\section{Lattice coordinates and exceptional fibers}\label{sec:lattice-boxes}

Membership in an approximation box of Lemma~\ref{lem:paired-boxes}
bounds linear forms such as $Mx_i-c_ix_0$, where $c_i/M$ is the
corresponding grid coordinate. We encode these bounds by lattices and
choose an adapted integral basis in each coordinate block. A box
requiring a larger coordinate bound has a shorter lattice vector;
such boxes are correspondingly fewer. This tradeoff will allow the
affine counting theorem to be summed over all boxes.
The use of a short vector to count the direction boxes in which it
can occur follows the method of Reuss
\cite[Sections~5--6]{Reuss2015PowerFree}. The common excess for the
two paired lattices below is tailored to our affine Hilbert bounds.

\subsection{A basis bound with uniform constants}

\begin{lemma}\label{lem:lattice-basis}
Let $\Lambda\subset\R^d$ be a full lattice of covolume $\Delta$,
and let $\lambda_1$ be its shortest nonzero Euclidean length.
There is a lattice basis $b_1,\ldots,b_d$ with
\[
 \prod_{i=1}^d\|b_i\|\le C_d\Delta.
\]
In this basis, every $v\in\Lambda$ with $\|v\|\le C$ has integral
coordinates of absolute value at most $CC_d/\lambda_1$.
\end{lemma}

\begin{proof}
We prove the product bound by induction on $d$. Choose a shortest
vector $b_1$, which is primitive, and project orthogonally to
$b_1^\perp$. The image $\Lambda'$ is a lattice of covolume
$\Delta/\|b_1\|$, and the projection kernel is $\Z b_1$.
Every nonzero vector of $\Lambda'$ has length at least
$\sqrt3\|b_1\|/2$: lift it and subtract an integer multiple of
$b_1$ until its parallel component has length at most $\|b_1\|/2$;
the resulting nonzero lattice vector still has length at least
$\|b_1\|$.

Apply induction to $\Lambda'$, and lift its basis vectors with
these reduced parallel components. Each lift has length at most
$2/\sqrt3$ times the projected length. Together with $b_1$ the
lifts form a basis of $\Lambda$, proving the asserted product
bound with a constant depending only on $d$. Finally, Cramer's
rule gives, for the coordinate $z_i$ of $v$,
\[
 |z_i|\le
 \frac{C\prod_{h\ne i}\|b_h\|}{\Delta}
 \le \frac{CC_d}{\|b_i\|}
 \le \frac{CC_d}{\lambda_1}.
\]
\end{proof}

\subsection{Grouping the paired boxes}

Fix a maximal-coordinate chart from Section~\ref{sec:arithmetic-boxes},
and reindex its coordinates as
$x=(x_0,\ldots,x_{d-1})$, $y=(y_0,\ldots,y_{d-1})$.
Thus $s_i=x_i/x_0$ and $t_i=y_i/y_0$ for $1\le i<d$.
Let $\eta_+,b_+>0$ be fixed upper exponents for which
\[
 |x_i|\ll X^{\eta_+/d},\qquad |y_i|\ll X^{b_+/d}.
\]
Choose fixed $w,\tau>0$ with
\begin{equation}\label{eq:lattice-parameters}
 m=(d-1)w<\min(\eta_+,b_+),\qquad
 u=\frac{\eta_+-m}{d},\qquad v=\frac{b_+-m}{d}.
\end{equation}
Use a prime $M\asymp X^w$ for the approximation grid. Such a prime
exists by Bertrand's postulate. All constants in this section are
uniform in the grid boxes and in $X$; they may depend on the fixed
parameters, the dimension, and the polynomial.

\begin{proposition}\label{prop:lattice-boxes}
For the paired boxes of Lemma~\ref{lem:paired-boxes}, assume the
coordinate bounds above and fix parameters satisfying
\eqref{eq:lattice-parameters}. The paired boxes can be grouped into
finitely many ranges indexed by $j=\ell\tau$, $\ell\in\Z_{\ge0}$,
with $j\le m/d$. The number of paired boxes in range $j$ is
\begin{equation}\label{eq:lattice-box-count}
 O\bigl(X^{m-dj}\bigr).
\end{equation}
For each paired box there are integral unimodular changes
$x\leftrightarrow x'$ and $y\leftrightarrow y'$ such that all
its selected points satisfy
\begin{equation}\label{eq:lattice-coordinate-bounds}
 |x'_i|\ll X^U,\qquad |y'_i|\ll X^V,
 \qquad U=u+j+\tau,\quad V=v+j+\tau.
\end{equation}
The matrices of these changes and their inverses have entries
bounded by a fixed power of $X$.
\end{proposition}

\begin{proof}
Suppose the $s$-box has grid base $(c_1/M,\ldots,c_{d-1}/M)$.
Its integers $c_i$ satisfy $|c_i|\ll M$. Consider the lattice
\[
 \Lambda_c=X^{-\eta_+/d}
 \left\{\bigl(q_0,Mq_1-c_1q_0,\ldots,Mq_{d-1}-c_{d-1}q_0\bigr):
 q\in\Z^d\right\}.
\]
The image of every counted $x$ has bounded Euclidean norm, since
\[
 Mx_i-c_ix_0=Mx_0(s_i-c_i/M),\qquad |s_i-c_i/M|\le M^{-1}.
\]
Let $U_*$ be the inverse shortest nonzero length of $\Lambda_c$.
Lemma~\ref{lem:lattice-basis} gives integral coordinates with
$|x'_i|\ll U_*$. The preimages of a lattice basis form a basis of
$\Z^d$, so this is a unimodular change. Apply the same construction
to the $t$-box, using the scaling $X^{-b_+/d}$, and denote its
inverse shortest length by $V_*$. Then $|y'_i|\ll V_*$.

Define the excess
\[
 J=\max\{1,U_*/X^u,V_*/X^v\}.
\]
Before the common scaling, every nonzero lattice vector is a
nonzero integer vector and has norm at least one. Consequently
\begin{equation}\label{eq:lattice-J-bound}
 U_*\le X^{\eta_+/d},\qquad V_*\le X^{b_+/d},
 \qquad 1\le J\le X^{m/d}.
\end{equation}
Put a pair in range $j$ when $X^j\le J<X^{j+\tau}$.
The exact upper bound in \eqref{eq:lattice-J-bound} shows that
only $0\le j\le m/d$ occur. In particular there are boundedly
many ranges and $X^{m/d-j}\ge1$ in every occupied range.
The upper bound for $J$ within its range proves
\eqref{eq:lattice-coordinate-bounds}.

For $j>0$, at least one of $U_*\ge X^{u+j}$ and
$V_*\ge X^{v+j}$ holds. In the first case a shortest lattice
vector gives a nonzero integer vector
\[
 q=(q_0,Mq'_1-c_1q_0,\ldots,Mq'_{d-1}-c_{d-1}q_0),
 \qquad \|q\|\le T:=X^{m/d-j}.
\]
Here $1\le T\le X^{m/d}$ and $T/M\ll X^{-w/d}$, so $T<M$
for sufficiently large $X$. If $q_0=0$, every other coordinate
would be divisible by $M$, contrary to $0<\|q\|<M$.
Thus $0<|q_0|<M$, and the primality of $M$ gives
\[
 c_i\equiv-q_i q_0^{-1}\pmod M\qquad(1\le i<d).
\]
Each vector $q$ determines the grid base up to boundedly many
possibilities, because $|c_i|\ll M$. There are $O(T^d)$ such
integer vectors. Hence there are $O(X^{m-dj})$ possible $s$-boxes
in this case. The $V_*$ case gives the identical bound for
$t$-boxes. Lemma~\ref{lem:paired-boxes} supplies boundedly many
partners in either direction, proving \eqref{eq:lattice-box-count}
for their union. When $j=0$, use the total bound
$O(M^{d-1})=O(X^m)$ from that lemma.

For completeness, the bases above also have polynomial height.
The covolume of $\Lambda_c$ is
$X^{-\eta_+}M^{d-1}\asymp X^{m-\eta_+}$, and its shortest length
is at least $X^{-\eta_+/d}$. The product bound in
Lemma~\ref{lem:lattice-basis} therefore gives
\[
 \max_i\|b_i\|\ll X^{m-\eta_+/d}.
\]
The inverse of the unscaled triangular map defining $\Lambda_c$
has bounded norm, since $|c_i|\ll M$. Its lattice basis thus
pulls back to integral vectors of size $O(X^m)$. Cofactors bound
the inverse unimodular matrix by a fixed power of $X$.
The argument for $y$ is the same.
\end{proof}

Linear changes preserve dimensions and degrees of all the cut
sets. They also preserve the weighted polynomial spaces within
each coordinate block: every $a_i$ is a linear combination of
the $x'_i$, which all have weight $U$, and every $b_i$ is a linear
combination of the $y'_i$, which all have weight $V$.
This assertion concerns the polynomial filtration, irrespective
of the coefficients of the linear forms. The numerical size
bounds used for integral counting are those on $x',y'$ in
\eqref{eq:lattice-coordinate-bounds}.

\subsection{Removing infinite fibers in the quartic case}

For $d=4$, a later Hilbert estimate requires $n$ to be algebraic
over the functions $a_1,\ldots,a_4$. We arrange this by removing
a small set of possible $x'$ from each initial cut set of
Proposition~\ref{prop:archimedean-cuts}.

\begin{lemma}\label{lem:quartic-fibers}
Suppose $d=4$, and let $Z$ be an initial bounded-degree set from
Proposition~\ref{prop:archimedean-cuts}, restricted to the open
$a_1a_2a_3a_4\ne0$. Use coordinates
$(a,n)$ there, with $b_i=(n-\theta_i)/a_i^2$.
Assume either that $\dim Z\le3$, or that $Z$ has the irreducible
equation $F(a,n)=0$ with positive degree in $n$.
The set of $a$ for which the $n$-fiber has dimension one has a
closure of dimension at most two and bounded degree.

In a box satisfying \eqref{eq:lattice-coordinate-bounds}, deleting
all selected points above this base costs $O_f(X^{2U})$ points.
On every irreducible subvariety of $Z$ meeting the remaining
selected set, $n$ is algebraic over the field generated by the
$a_i$.
\end{lemma}

\begin{proof}
Substitute $b_i=(n-\theta_i)/a_i^2$ into bounded-degree defining
equations for $Z$ and clear their denominators. On the specified
open, a fiber has dimension one precisely when each resulting
polynomial in $n$ vanishes identically. Thus its base is the common
zero set of their coefficient polynomials, restricted to the
$a$-torus. These equations have bounded degrees; retaining their
relevant components and taking closure preserves a bounded degree.
Components contained entirely in $a_1a_2a_3a_4=0$ are not included.

If $\dim Z\le3$, the fiber-dimension inequality gives dimension
at most two for this base. In the other case the coefficients of
$F$ as a polynomial in $n$ have no common nonconstant factor:
such a factor would factor the irreducible polynomial $F$, whose
$n$-degree is positive. Their common zero set therefore has
codimension at least two in four-dimensional $a$-space. Indeed,
a codimension-one component would give an irreducible polynomial
dividing every coefficient. This proves the dimension bound in
both cases.

The invertible linear change from $a$ to $x'$ preserves this
dimension and degree. A bounded-degree algebraic set of dimension
at most two contains $O(B^2)$ integer points in a box of radius
$B\ge1$, with a constant independent of its coefficients.
To recall the elementary argument, split into irreducible
components and slice a positive-dimensional component by a
nonconstant coordinate. Each integer slice is a proper
hyperplane section of smaller dimension and bounded degree;
induction ends with the degree bound for a zero-dimensional set.
Taking $B\asymp X^U$ bounds the exceptional $x'$ by $O(X^{2U})$.
Each such $x'$ determines $x$, and the fixed-$x$ conclusion of
Lemma~\ref{lem:balanced-factorization} gives only $O_f(1)$ selected
points above it.

Finally, suppose $Y\subseteq Z$ is irreducible and $n$ is
transcendental over the field generated by its $a_i$. Every
polynomial relation in $n$ from the substituted equations then
has all coefficients zero on $Y$. Its entire $a$-image lies in
the base just removed. Consequently such a $Y$ meets no remaining
selected point, proving the last assertion.
\end{proof}

Proposition~\ref{prop:lattice-boxes} supplies the integral boxes
for the affine counting theorem, and Lemma~\ref{lem:quartic-fibers}
supplies the additional projection property needed in the quartic
case. The next sections establish the Hilbert bounds on the sets
inside these boxes; the final summation will combine their cost
with the factor $X^{m-dj}$ in \eqref{eq:lattice-box-count}.
